% questions.tex —— 2013 年·江苏卷 · 题目内容片段（无答案）
% 由 tools/generate_projects.py 依来源总账自 2013/2013_数学真题_江苏_试题.md 生成，勿手改（改源 MD 后重跑）


\begin{gknotice}
\textbf{来源说明}：本题面依据固定源 PDF 的两页 A3 题面逐题转录，原卷印刷页码为 815—816。第 815 面为填空题 1—9，第 816 面为填空题 10—14、解答题 15—20、附加题 21—23。第 21 题为四选二，A、B、C、D 四个小题全部收录；为便于检索，四个小题分别以（1）—（4）标记，原卷的【A】【B】【C】【D】标识一并保留。

\textbf{解析归属}：固定源不附答案或解析；配套答案与解析均由本库按本题面独立复算。

\textbf{图形说明}：第 5 题算法流程图、第 8 题三棱柱、第 16 题三棱锥、第 17 题坐标系与直线、第 18 题 A-B-C 路线、第 21A 题圆与切线、第 22 题直三棱柱，均由项目内 TikZ 矢量源重绘。
\end{gknotice}

\gkpart{一、填空题}

\begin{question}
\gkti{1}函数 $y=3\sin\left(2x+\dfrac{\pi}{4}\right)$ 的最小正周期为\gkblank{}。

\end{question}
\begin{question}
\gkti{2}设 $z=(2-\mathrm i)^2$（$\mathrm i$ 为虚数单位），则复数 $z$ 的模为\gkblank{}。

\end{question}
\begin{question}
\gkti{3}双曲线 $\dfrac{x^2}{16}-\dfrac{y^2}{9}=1$ 的两条渐近线的方程为\gkblank{}。

\end{question}
\begin{question}
\gkti{4}集合 $\{-1,0,1\}$ 共有\gkblank{}个子集。

\Needspace{0.4\textheight}
\end{question}
\begin{question}
\gkti{5}如图是一个算法的流程图，则输出的 $n$ 的值是\gkblank{}。
\nopagebreak[4]
\begin{center}% ============================================================
% f2013_jiangsu_q05_flowchart.tex
% 2013 江苏卷 第 5 题 算法流程图（TikZ 矢量重绘）
% n←1，a←2；a<20 时 a←3a+2、n←n+1，否则输出 n
% ============================================================
\begin{tikzpicture}[
  >=Stealth, line width=0.6pt, font=\small,
  box/.style={draw, rounded corners=3pt, minimum width=2.7cm,
    minimum height=0.85cm, align=center},
  proc/.style={draw, minimum width=2.7cm, minimum height=0.85cm, align=center},
  dec/.style={draw, diamond, aspect=1.9, minimum width=3.4cm,
    minimum height=1.35cm, inner sep=1pt, align=center},
  io/.style={draw, trapezium, trapezium left angle=70,
    trapezium right angle=110, minimum width=2.9cm,
    minimum height=0.85cm, align=center}
]

  \node[box]  (start) {开始};
  \node[proc, below=7mm of start] (none) {$n\leftarrow 1$};
  \node[proc, below=6mm of none]  (atwo) {$a\leftarrow 2$};
  \node[dec,  below=8mm of atwo]  (test) {$a<20$?};
  \node[proc, right=17mm of test] (athree) {$a\leftarrow 3a+2$};
  \node[proc, above=10mm of athree] (inc) {$n\leftarrow n+1$};
  \node[io,   below=8mm of test]  (out) {输出 $n$};
  \node[box,  below=7mm of out]   (end) {结束};

  \draw[->] (start) -- (none);
  \draw[->] (none) -- (atwo);
  \draw[->] (atwo) -- (test);
  \coordinate (join) at ($(atwo.south)!0.55!(test.north)$);
  \draw[->] (test.east) -- node[above] {$Y$} (athree.west);
  \draw[->] (athree.north) -- (inc.south);
  \draw (inc.west) -- ++(-0.25,0) |- (join);
  \draw[->] (test.south) -- node[right] {$N$} (out.north);
  \draw[->] (out.south) -- (end.north);

\end{tikzpicture}
\end{center}
\nopagebreak[4]
\end{question}
\begin{question}
\gkti{6}抽样统计甲、乙两位射击运动员的 $5$ 次训练成绩（单位：环），结果如下：

\begin{center}
\begin{tabular}{llllll}
\toprule
运动员 & 第 1 次 & 第 2 次 & 第 3 次 & 第 4 次 & 第 5 次 \\
\midrule
甲 & 87 & 91 & 90 & 89 & 93 \\
乙 & 89 & 90 & 91 & 88 & 92 \\
\bottomrule
\end{tabular}
\end{center}

则成绩较为稳定（方差较小）的那位运动员成绩的方差为\gkblank{}。

\end{question}
\begin{question}
\gkti{7}现有某类病毒记作 $X_mY_n$，其中正整数 $m$，$n$（$m\leqslant7$，$n\leqslant9$）可以任意选取，则 $m$，$n$ 都取到奇数的概率为\gkblank{}。

\Needspace{0.4\textheight}
\end{question}
\begin{question}
\gkti{8}如图，在三棱柱 $A_1B_1C_1-ABC$ 中，$D$，$E$，$F$ 分别是 $AB$，$AC$，$AA_1$ 的中点。设三棱锥 $F-ADE$ 的体积为 $V_1$，三棱柱 $A_1B_1C_1-ABC$ 的体积为 $V_2$，则 $V_1:V_2=$\gkblank{}。
\nopagebreak[4]
\begin{center}% ============================================================
% f2013_jiangsu_q08_prism.tex
% 2013 江苏卷 第 8 题 三棱柱 A1B1C1-ABC 与截面 F-ADE
% D、E、F 分别为 AB、AC、AA1 的中点
% ============================================================
\begin{tikzpicture}[
  line width=0.6pt, font=\small, >=Stealth
]

  \coordinate (A)  at (0.2,0);
  \coordinate (B)  at (4.4,0.35);
  \coordinate (C)  at (2.0,1.95);
  \coordinate (A1) at (0.2,3.55);
  \coordinate (B1) at (4.4,3.90);
  \coordinate (C1) at (2.0,5.50);
  \coordinate (D)  at ($(A)!0.5!(B)$);
  \coordinate (E)  at ($(A)!0.5!(C)$);
  \coordinate (F)  at ($(A)!0.5!(A1)$);

  % 底面（AB、BC 可见，AC 为被遮住的边）
  \draw (A) -- (B) -- (C);
  \draw[dashed] (A) -- (C);

  % 侧棱
  \draw (A) -- (A1);
  \draw (B) -- (B1);
  \draw[dashed] (C) -- (C1);

  % 顶面
  \draw (A1) -- (B1) -- (C1) -- cycle;

  % 截面 F-ADE
  \draw (F) -- (D) -- (E);
  \draw[dashed] (E) -- (F);

  % 标注
  \node[below left]  at (A)  {$A$};
  \node[below right] at (B)  {$B$};
  \node[right]       at (C)  {$C$};
  \node[left]        at (A1) {$A_1$};
  \node[right]       at (B1) {$B_1$};
  \node[above]       at (C1) {$C_1$};
  \node[below]       at (D)  {$D$};
  \node[right]       at (E)  {$E$};
  \node[left]        at (F)  {$F$};

\end{tikzpicture}
\end{center}
\nopagebreak[4]
\end{question}
\begin{question}
\gkti{9}抛物线 $y=x^2$ 在 $x=1$ 处的切线与两坐标轴围成的三角形区域为 $D$（包含三角形内部与边界）。若点 $P(x,y)$ 是区域 $D$ 内的任意一点，则 $x+2y$ 的取值范围是\gkblank{}。

\end{question}
\begin{question}
\gkti{10}设 $D$，$E$ 分别是 $\triangle ABC$ 的边 $AB$，$BC$ 上的点，$AD=\dfrac12AB$，$BE=\dfrac23BC$。若 $\overrightarrow{DE}=\lambda_1\overrightarrow{AB}+\lambda_2\overrightarrow{AC}$（$\lambda_1$，$\lambda_2$ 为实数），则 $\lambda_1+\lambda_2$ 的值为\gkblank{}。

\end{question}
\begin{question}
\gkti{11}已知 $f(x)$ 是定义在 $\mathbf R$ 上的奇函数。当 $x>0$ 时，$f(x)=x^2-4x$，则不等式 $f(x)>x$ 的解集用区间表示为\gkblank{}。

\end{question}
\begin{question}
\gkti{12}在平面直角坐标系 $xOy$ 中，椭圆 $C$ 的标准方程为 $\dfrac{x^2}{a^2}+\dfrac{y^2}{b^2}=1$（$a>b>0$），右焦点为 $F$，右准线为 $l$，短轴的一个端点为 $B$。设原点到直线 $BF$ 的距离为 $d_1$，$F$ 到 $l$ 的距离为 $d_2$。若 $d_2=\sqrt6\,d_1$，则椭圆 $C$ 的离心率为\gkblank{}。

\end{question}
\begin{question}
\gkti{13}在平面直角坐标系 $xOy$ 中，设定点 $A(a,a)$，$P$ 是函数 $y=\dfrac1x$（$x>0$）图象上一动点。若点 $P$，$A$ 之间的最短距离为 $2\sqrt2$，则满足条件的实数 $a$ 的所有值为\gkblank{}。

\end{question}
\begin{question}
\gkti{14}在正项等比数列 $\{a_n\}$ 中，$a_5=\dfrac12$，$a_6+a_7=3$，则满足 $a_1+a_2+\cdots+a_n>a_1a_2\cdots a_n$ 的最大正整数 $n$ 的值为\gkblank{}。

\end{question}
\gkpart{二、解答题}

\begin{question}
\gkti{15}已知 $\boldsymbol a=(\cos\alpha,\sin\alpha)$，$\boldsymbol b=(\cos\beta,\sin\beta)$，$0<\beta<\alpha<\pi$。

（1） 若 $|\boldsymbol a-\boldsymbol b|=\sqrt2$，求证：$\boldsymbol a\perp\boldsymbol b$；

（2） 设 $\boldsymbol c=(0,1)$，若 $\boldsymbol a+\boldsymbol b=\boldsymbol c$，求 $\alpha$，$\beta$ 的值。

\Needspace{0.4\textheight}
\end{question}
\begin{question}
\gkti{16}如图，在三棱锥 $S-ABC$ 中，平面 $SAB\perp$ 平面 $SBC$，$AB\perp BC$，$AS=AB$。过 $A$ 作 $AF\perp SB$，垂足为 $F$，点 $E$，$G$ 分别是棱 $SA$，$SC$ 的中点。求证：

（1） 平面 $EFG\parallel$ 平面 $ABC$；

（2） $BC\perp SA$。
\nopagebreak[4]
\begin{center}% ============================================================
% f2013_jiangsu_q16_tetrahedron.tex
% 2013 江苏卷 第 16 题 三棱锥 S-ABC
% E、G 为 SA、SC 中点，F 为 SB 上一点（AF⊥SB）
% ============================================================
\begin{tikzpicture}[
  line width=0.6pt, font=\small, >=Stealth
]

  \coordinate (S) at (1.8,4.3);
  \coordinate (A) at (0,1.4);
  \coordinate (B) at (2.0,-0.3);
  \coordinate (C) at (4.9,1.0);
  \coordinate (E) at ($(S)!0.5!(A)$);
  \coordinate (G) at ($(S)!0.5!(C)$);
  \coordinate (F) at ($(S)!0.5!(B)$);

  % 棱
  \draw (S) -- (A) -- (B) -- (C) -- (S);
  \draw (S) -- (B);
  \draw[dashed] (A) -- (C);

  % AF 与截面 EFG
  \draw (A) -- (F);
  \draw (E) -- (F) -- (G);
  \draw[dashed] (E) -- (G);

  % 标注
  \node[above]      at (S) {$S$};
  \node[left]       at (A) {$A$};
  \node[below]      at (B) {$B$};
  \node[right]      at (C) {$C$};
  \node[left]       at (E) {$E$};
  \node[right]      at (G) {$G$};
  \node[below right] at (F) {$F$};

\end{tikzpicture}
\end{center}
\nopagebreak[4]
\Needspace{0.4\textheight}
\end{question}
\begin{question}
\gkti{17}如图，在平面直角坐标系 $xOy$ 中，点 $A(0,3)$，直线 $l:y=2x-4$。设圆 $C$ 的半径为 $1$，圆心在 $l$ 上。

（1） 若圆心 $C$ 也在直线 $y=x-1$ 上，过点 $A$ 作圆 $C$ 的切线，求切线的方程；

（2） 若圆 $C$ 上存在点 $M$，使 $MA=2MO$，求圆心 $C$ 的横坐标 $a$ 的取值范围。
\nopagebreak[4]
\begin{center}% ============================================================
% f2013_jiangsu_q17_line.tex
% 2013 江苏卷 第 17 题 坐标系、点 A(0,3) 与直线 l: y=2x-4
% 原卷图仅画坐标轴、点 A 和直线 l（圆 C 未画出）
% ============================================================
\begin{tikzpicture}[
  x=0.72cm, y=0.62cm, line width=0.6pt, font=\small, >=Stealth
]

  % 坐标轴
  \draw[->] (-0.8,0) -- (4.7,0) node[below] {$x$};
  \draw[->] (0,-4.9) -- (0,4.6) node[left] {$y$};
  \node[below left] at (0,0) {$O$};

  % 刻度
  \foreach \x in {1,2,3,4} \draw (\x,0.1) -- (\x,-0.1);
  \foreach \y in {-4,-3,-2,-1,1,2,3,4} \draw (0.1,\y) -- (-0.1,\y);

  % 点 A
  \fill (0,3) circle (1.6pt);
  \node[left] at (0,3) {$A$};

  % 直线 l: y=2x-4
  \draw[thick] (-0.45,-4.9) -- (4.45,4.9);
  \node[right] at (3.9,2.9) {$l$};

\end{tikzpicture}
\end{center}
\nopagebreak[4]
\Needspace{0.4\textheight}
\end{question}
\begin{question}
\gkti{18}如图，游客从某旅游景区的景点 $A$ 处下山至 $C$ 处有两种路径。一种是从 $A$ 沿直线步行到 $C$，另一种是先从 $A$ 沿索道乘缆车到 $B$，然后从 $B$ 沿直线步行到 $C$。现有甲、乙两位游客从 $A$ 处下山，甲沿 $AC$ 匀速步行，速度为 $50\ \mathrm{m/min}$。在甲出发 $2\ \mathrm{min}$ 后，乙从 $A$ 乘缆车到 $B$，在 $B$ 处停留 $1\ \mathrm{min}$ 后，再从 $B$ 匀速步行到 $C$。假设缆车匀速直线运动的速度为 $130\ \mathrm{m/min}$，山路 $AC$ 长为 $1260\ \mathrm{m}$，经测量，$\cos A=\dfrac{12}{13}$，$\cos C=\dfrac35$。

（1） 求索道 $AB$ 的长；

（2） 问乙出发多少分钟后，乙在缆车上与甲的距离最短？

（3） 为使两位游客在 $C$ 处互相等待的时间不超过 $3$ 分钟，乙步行的速度应控制在什么范围内？
\nopagebreak[4]
\begin{center}% ============================================================
% f2013_jiangsu_q18_route.tex
% 2013 江苏卷 第 18 题 A-B-C 两条路径示意
% A 沿 AC 步行；A 乘缆车到 B，再步行到 C
% ============================================================
\begin{tikzpicture}[
  line width=0.6pt, font=\small
]

  \coordinate (A) at (6.6,4.2);
  \coordinate (B) at (3.4,3.1);
  \coordinate (C) at (0.4,0.3);

  \draw (A) -- (B) -- (C);
  \draw (A) -- (C);

  \node[right] at (A) {$A$};
  \node[above] at (B) {$B$};
  \node[below left] at (C) {$C$};

\end{tikzpicture}
\end{center}
\nopagebreak[4]
\end{question}
\begin{question}
\gkti{19}设 $\{a_n\}$ 是首项为 $a$，公差为 $d$ 的等差数列（$d\neq0$），$S_n$ 是其前 $n$ 项的和。记 $b_n=\dfrac{nS_n}{n^2+c}$，$n\in\mathbf N^*$，其中 $c$ 为实数。

（1） 若 $c=0$，且 $b_1$，$b_2$，$b_4$ 成等比数列，证明：$S_{nk}=n^2S_k$（$k,n\in\mathbf N^*$）；

（2） 若 $\{b_n\}$ 是等差数列，证明：$c=0$。

\end{question}
\begin{question}
\gkti{20}设函数 $f(x)=\ln x-ax$，$g(x)=\mathrm e^x-ax$，其中 $a$ 为实数。

（1） 若 $f(x)$ 在 $(1,+\infty)$ 上是单调减函数，且 $g(x)$ 在 $(1,+\infty)$ 上有最小值，求 $a$ 的取值范围；

（2） 若 $g(x)$ 在 $(-1,+\infty)$ 上是单调增函数，试求 $f(x)$ 的零点个数，并证明你的结论。

\end{question}
\gkpart{三、附加题}

\Needspace{0.4\textheight}
\begin{question}
\gkti{21}四选二。

（1） 【A】如图，$AB$ 和 $BC$ 分别与圆 $O$ 相切于点 $D$，$C$，$AC$ 经过圆心 $O$，且 $BC=2OC$。求证：$AC=2AD$。
\nopagebreak[4]
\begin{center}% ============================================================
% f2013_jiangsu_q21a_circle.tex
% 2013 江苏卷 第 21A 题 圆 O 与切线 AB、BC
% AB、BC 分别切圆 O 于 D、C，AC 经过圆心 O，BC=2OC
% ============================================================
\begin{tikzpicture}[
  x=1.05cm, y=1.05cm, line width=0.6pt, font=\small, >=Stealth
]

  \coordinate (O) at (0,0);
  \coordinate (C) at (-1.5,0);
  \coordinate (A) at (2.5,0);
  \coordinate (B) at (-1.5,3.0);
  \coordinate (D) at (0.9,1.2);

  % 圆
  \draw (O) circle (1.5);

  % 切线与割线
  \draw (C) -- (B);
  \draw (C) -- (A);
  \draw (B) -- (A);

  % 圆心与切点
  \fill (O) circle (1.5pt);
  \fill (D) circle (1.5pt);
  \fill (C) circle (1.2pt);

  % C 处直角标记
  \draw (C) ++(0,0.28) -- ++(0.28,0) -- ++(0,-0.28);

  % 标注
  \node[below left] at (C) {$C$};
  \node[right] at (A) {$A$};
  \node[above] at (B) {$B$};
  \node[above right] at (D) {$D$};
  \node[below] at (O) {$O$};

\end{tikzpicture}
\end{center}
\nopagebreak[4]
（2） 【B】已知矩阵 $A=\begin{bmatrix}-1&0\\0&2\end{bmatrix}$，$B=\begin{bmatrix}1&2\\0&6\end{bmatrix}$，求矩阵 $A^{-1}B$。

（3） 【C】在平面直角坐标系 $xOy$ 中，直线 $l$ 的参数方程为 $\begin{cases}x=t+1,\\y=2t,\end{cases}$（$t$ 为参数），曲线 $C$ 的参数方程为 $\begin{cases}x=2\tan^2\theta,\\y=2\tan\theta,\end{cases}$（$\theta$ 为参数）。试求直线 $l$ 和曲线 $C$ 的普通方程，并求出它们的公共点的坐标。

（4） 【D】已知 $a\geqslant b>0$，求证：$2a^3-b^3\geqslant2ab^2-a^2b$。

\Needspace{0.4\textheight}
\end{question}
\begin{question}
\gkti{22}如图，在直三棱柱 $A_1B_1C_1-ABC$ 中，$AB\perp AC$，$AB=AC=2$，$A_1A=4$，点 $D$ 是 $BC$ 的中点。

（1） 求异面直线 $A_1B$ 与 $C_1D$ 所成角的余弦值；

（2） 求平面 $ADC_1$ 与平面 $ABA_1$ 所成二面角的正弦值。
\nopagebreak[4]
\begin{center}% ============================================================
% f2013_jiangsu_q22_prism.tex
% 2013 江苏卷 第 22 题 直三棱柱 A1B1C1-ABC
% AB⊥AC，AB=AC=2，A1A=4，D 为 BC 中点
% ============================================================
\begin{tikzpicture}[
  line width=0.6pt, font=\small, >=Stealth
]

  \coordinate (A)  at (1.3,0.85);
  \coordinate (B)  at (0,0);
  \coordinate (C)  at (3.4,0);
  \coordinate (A1) at (1.3,3.9);
  \coordinate (B1) at (0,3.05);
  \coordinate (C1) at (3.4,3.05);
  \coordinate (D)  at ($(B)!0.5!(C)$);

  % 底面
  \draw (A) -- (B) -- (C);
  \draw[dashed] (A) -- (C);

  % 侧棱
  \draw (B) -- (B1);
  \draw (C) -- (C1);
  \draw[dashed] (A) -- (A1);

  % 顶面
  \draw (A1) -- (B1) -- (C1) -- cycle;

  % 题设中的线段
  \draw[dashed] (A1) -- (B);
  \draw (C1) -- (D);
  \draw[dashed] (A) -- (D);
  \draw[dashed] (A) -- (C1);

  % 标注
  \node[above left]  at (A)  {$A$};
  \node[below left] at (B)  {$B$};
  \node[right]      at (C)  {$C$};
  \node[left]       at (A1) {$A_1$};
  \node[left]       at (B1) {$B_1$};
  \node[right]      at (C1) {$C_1$};
  \node[below]      at (D)  {$D$};

\end{tikzpicture}
\end{center}
\nopagebreak[4]
\end{question}
\begin{question}
\gkti{23}设数列 $\{a_n\}$：$1,-2,-2,3,3,3,-4,-4,-4,-4,\cdots$，
\nopagebreak[4]
\[\underbrace{(-1)^{k-1}k,\cdots,(-1)^{k-1}k}_{k\text{ \text{个}}},\cdots,\]
即当 $\dfrac{(k-1)k}{2}<n\leqslant\dfrac{k(k+1)}{2}$（$k\in\mathbf N^*$）时，$a_n=(-1)^{k-1}k$。记 $S_n=a_1+a_2+\cdots+a_n$（$n\in\mathbf N^*$）。对于 $l\in\mathbf N^*$，定义集合 $P_l=\{n\mid S_n\text{ \text{是} }a_n\text{ \text{的整数倍}},\,n\in\mathbf N^*,\,\text{\text{且} }1\leqslant n\leqslant l\}$。

（1） 求集合 $P_{11}$ 中元素的个数；

（2） 求集合 $P_{2000}$ 中元素的个数。

\end{question}